% !TEX program = xelatex
% !TEX options = -synctex=1 -interaction=nonstopmode -file-line-error
% example-poster-tikz.tex
% Demonstrates rossposter (tikzposter-based): title block, content
% blocks, a figure, QR code, and compact references.

\begin{filecontents}{example-poster.bib}
@article{smith2024sleep,
  author  = {Smith, J. A. and Lee, K. R.},
  title   = {Sleep architecture and cognitive decline},
  journal = {Sleep},
  year    = {2024},
  volume  = {47},
  pages   = {1--12},
}
@article{randolph1998rbans,
  author  = {Randolph, C.},
  title   = {RBANS manual},
  journal = {Psychological Corporation},
  year    = {1998},
}
\end{filecontents}

\documentclass[palette=uom, size=a0]{rossposter}
\addbibresource{example-poster.bib}

% ── Title and authors ────────────────────────────────────────────
\posterauthor{Ross A.\ Dunne}
  {University of Limerick}{ross.dunne@ul.ie}
\posterauthor{Sarah K.\ Mitchell}
  {Trinity College Dublin}{s.mitchell@tcd.ie}
\posterauthor{Ann-Marie Moriarty}
  {University of Limerick}{ann-marie.moriarty@ul.ie}

\begin{document}

\postertitle{Slow-Wave Sleep Predicts Delayed Memory via \\ Depressive
             Symptoms in Community-Dwelling \\ Older Adults}

% ── Left column ──────────────────────────────────────────────────
\begin{columns}
\column{0.33}

\posterblock{Background}{%
  Sleep disturbance is an established risk factor for cognitive decline.
  Slow-wave sleep (SWS) supports memory consolidation, yet the mediating
  role of affective symptoms is unclear \cite{smith2024sleep}.

  \textbf{Aims:} (1) Examine the direct effect of SWS\% on RBANS
  delayed memory; (2) Test depressive and anxiety symptoms as parallel
  mediators.
}

\posterblock{Methods}{%
  \begin{itemize}
    \item $N = 312$ adults aged 60--85 (58\% female)
    \item Single-night in-lab polysomnography
    \item DASS-21 (depression and anxiety subscales)
    \item RBANS \cite{randolph1998rbans} five index scores
    \item Parallel mediation with 5,000 bootstrap resamples
  \end{itemize}
}

% ── Middle column ────────────────────────────────────────────────
\column{0.34}

\posterblock{Results}{%
  SWS\% significantly predicted delayed memory
  (\effectsize{\beta}{$-$0.24}, \pval{0.001}).

  The indirect effect through depressive symptoms was significant
  (\ci{$-$0.18}{$-$0.34}{$-$0.08}, \pval{0.002}).
  No indirect effect through anxiety (\pval{0.41}).

  \vspace{8mm}
  \posterfigure[0.85\linewidth]{example-image-a}{%
    Mediation model: standardised path coefficients.}
  \figurenote{Dashed path = non-significant indirect effect through
    anxiety symptoms.}
}

% ── Right column ─────────────────────────────────────────────────
\column{0.33}

\posterblock{Conclusions}{%
  \begin{itemize}
    \item Depressive symptoms partially mediate the SWS--memory link.
    \item Targeting mood in older adults with sleep complaints may
          yield cognitive benefits.
    \item Longitudinal replication with multi-night PSG is needed.
  \end{itemize}
}

\posterblock{Contact}{%
  \centering
  \posterqr{https://osf.io/abcde}{Scan for data \& materials}
}

\posterreferences

\end{columns}

% ── Footer ───────────────────────────────────────────────────────
\posterfooter
  {Department of Psychology, University of Limerick}
  {INS 53rd Annual Meeting --- New York, February 2026}
  {Funded by Irish Research Council (GOIPG/2024/1234)}

\end{document}
